\subsection{Coordinated screens and reciprocal-radius duality}
\label{exc:pantallas-dualidad}

The reduction selected by \(K\) and the Leech construction
arise from different readings of the same incidence. To specify
their coordination, the algebraic maps preserving
both data are brought together here. The rank-eleven space and rank-twenty-four lattice
are not identified; their respective quotients and the
involution acting on the route-and-memory factor are specified.

\subsubsection{Residue, quotient, and change of section}

The lifted arithmetic evaluation of APP preserves
\(n=9Q+r\), with \(r\in\{1,\ldots,9\}\) in the positive
residual section. TPK transport preserves these coordinates
separately, together with their orientation and the memory of their changes. To
represent an exchange between the two coordinates, a domain stable
under exchange is required. The ambient space
\(\ZZ^2\) is therefore used, whose signed integers prolong the coordinates of
the chart; this does not assert that all its pairs are admissible
trajectories of the same TPK history.

On \(\mathscr D_0=\ZZ^2\times\RR_{>0}\), the declared quadratic
realization is
\begin{equation}
 q_R(m,w)=\frac{m^2}{R^2}+w^2R^2,
 \qquad \mathsf T_0(m,w,R)=(w,m,R^{-1}).
 \label{eq:dualidad-normalizada-k}
\end{equation}
The variable \(R\) is a dimensionless scale of the reading chart.
The symbols \(m,w\) here designate the two transported integer
coordinates. Their subsequent recognition as momentum and winding
maintains that distinction between construction and realization.

\begin{proposition}[Exchange and quadratic form]
\label{prop:dualidad-radios-k}
The map \(\mathsf T_0\) is an involution and satisfies
\[
 q_{R^{-1}}(w,m)=q_R(m,w).
\]
On \(\ell^2(\ZZ^2)\), the unitary operator
\(U_T|m,w\rangle=|w,m\rangle\) intertwines the diagonal
operators \(H(R)|m,w\rangle=q_R(m,w)|m,w\rangle\):
\[
 U_TH(R)U_T^{-1}=H(R^{-1}).
\]
\end{proposition}
\begin{proof}
Two exchanges recover the coordinates, and two inversions recover
\(R\). Substitution gives
\(q_{R^{-1}}(w,m)=w^2R^2+m^2/R^2\). Exchange permutes
the orthonormal basis and is therefore unitary, and that equality verifies the
intertwining on each basis vector. It extends to the
domains of the diagonal operators through the same equality of weights.
\end{proof}

The residual section matters. For \(n=9\), the pairs
\((r,Q)=(9,0)\) and \((0,1)\) represent the same integer;
their evaluations \(81/R^2\) and \(R^2\) do not coincide for every
\(R\). Change of section is incorporated exactly through
\[
 L_9=\ZZ(9,-1),\qquad S(x_1,x_2)=(x_2,x_1).
\]
On quotient classes we define
\begin{equation}
 \mathcal E_R^{L_9}([x])=\min_{\ell\in L_9}q_R(x+\ell),
 \qquad
 \mathsf T_{\mathrm{cov}}([x]_{L_9},R)
       =([Sx]_{SL_9},R^{-1}).
 \label{eq:dualidad-cociente-k}
\end{equation}

\begin{proposition}[Duality compatible with change of section]
The two expressions in \eqref{eq:dualidad-cociente-k} are well
defined. The positive and zero sections, with the corresponding
quotient increment, represent the same class. Moreover,
\[
 \mathcal E_{R^{-1}}^{SL_9}([Sx])=\mathcal E_R^{L_9}([x]).
\]
\end{proposition}
\begin{proof}
Replacing \(x\) by \(x+\ell_0\) permutes the candidates for the minimum.
The minimum exists because \(q_R(x+k(9,-1))\) is a coercive quadratic
in \(k\in\ZZ\). The change \((9,Q)\mapsto(0,Q+1)\)
has difference \((9,-1)\). Finally,
\[
 \min_{\ell\in L_9}q_{R^{-1}}(Sx+S\ell)
 =\min_{\ell\in L_9}q_R(x+\ell).
\]
Since \(S^2=I\), the inverse exchange recovers the class and scale.
\end{proof}

The equality depends on also transporting the equivalence relation,
from \(L_9\) to \(SL_9\). This preserves the change-of-section
information that a mere substitution of representative \(9\) by
\(0\) would lose.

\subsubsection{The lattice screen and its coordination with \(K\)}

Let \(\Lambda=N_\alpha\), the Leech lattice constructed in
\ref{exc:leech}. Add the integral hyperbolic plane
\(U=\ZZ e_+\oplus\ZZ e_-\), with
\(e_+^2=e_-^2=0\) and \(\langle e_+,e_-\rangle=1\).
Then \(L_{26}=\Lambda\oplus U\) has signature \((25,1)\), and
\[
 e_+^\perp=\Lambda\oplus\ZZ e_+,
 \qquad e_+^\perp/\ZZ e_+\cong\Lambda.
\]
Indeed, if \(y=\lambda+ae_++be_-\), the condition
\(\langle y,e_+\rangle=0\) is equivalent to \(b=0\), and the quotient
eliminates exactly the summand \(ae_+\). Its map is
\(q_{24}(\lambda+ae_+)=\lambda\).

For the reduction \(P_{10}\) already proved, let
\[
 \mathscr G_K=\{(v_{11},P_{10}v_{11}):v_{11}\in V_{11}\}.
\]
Define
\begin{equation}
 \mathcal R_K(v,d,y)
 =\bigl((P_{11}v,P_{10}v),d,q_{24}(y)\bigr),
 \label{eq:k-mapa-pantallas}
\end{equation}
with domain \(\RR^{12}\times\mathscr D_0\times e_+^\perp\)
and codomain \(\mathscr G_K\times\mathscr D_0\times\Lambda\).

\begin{theorem}[Isomorphism of the quotient screens]
\label{thm:k-pantallas-coordinadas}
The map \eqref{eq:k-mapa-pantallas} induces an isomorphism
\[
 (\RR^{12}/\RR\mathbf1)\times\mathscr D_0
       \times(e_+^\perp/\ZZ e_+)
 \ \cong\ \mathscr G_K\times\mathscr D_0\times\Lambda.
\]
It intertwines the transformations acting by \(\mathsf T_0\)
on the factor \(\mathscr D_0\) and by the identity on the others.
\end{theorem}
\begin{proof}
\(P_{11}\) induces an isomorphism from \(\RR^{12}/\RR\mathbf1\)
onto \(V_{11}\). Since \(P_{10}P_{11}=P_{10}\), the map
\([v]\mapsto(P_{11}v,P_{10}v)\) identifies that quotient with
the graph \(\mathscr G_K\); the inverse takes the first component
and its class. The lattice quotient has just been identified through
\(q_{24}\). The remaining factor is preserved by the identity.
The product of the three maps proves the isomorphism. The
involution and projectors act on distinct factors, so
their compositions commute. The quadratic form is preserved
by Proposition \ref{prop:dualidad-radios-k}.
\end{proof}

The graph retains \(v_{11}\) together with its projection. This
isomorphism therefore does not assert that \(V_{11}\) and \(V_{10}\) are isomorphic:
the isolated reduction has kernel \(\ell_K\), whereas the graph
preserves the preceding coordinate. This specification identifies what each
map preserves and allows the screens \(11\to10\) and
\(26\to24\) to be linked without confusing them.

Taken together, \(K\) determines the reduction direction; the dodecaphase
flag determines the lattice neighbor used for Moonshine;
and exchange of the route and memory coordinates preserves the
quadratic form under reciprocal radii. Application to a physical
string theory or to an eleven-dimensional space additionally requires
its action, fields, and realization operators. The result brought together
here is the algebraic map with its proofs, constituting an
explicit antecedent for that subsequent development.
